\documentclass[11pt,a4paper]{article}
\usepackage[utf8]{inputenc}
\usepackage[T1]{fontenc}
\usepackage{amsmath,amssymb}
\usepackage{booktabs}
\usepackage{geometry}
\geometry{left=2.2cm,right=2.2cm,top=2.2cm,bottom=2.2cm}
\usepackage{url}
\usepackage{hyperref}
\hypersetup{colorlinks=true,linkcolor=blue,citecolor=blue,urlcolor=blue}
\usepackage{enumitem}

\title{\textbf{CROSS: Cross Router Of Selector System}\\
\large Program-Level MoE Fusion over a Frozen Language Model\\
\large --- A Positive Fusion Result and a Negative Selector Result}
\author{Weikai Yue, Hy 4 flash}
\date{Independent research report --- Draft v0.1 --- 2026-11-26}

\begin{document}
\maketitle

\begin{center}
\textit{Follow-up to Amoeba and BRICK \quad $\cdot$ \quad CPU-only empirical study \quad $\cdot$ \quad All results, including negative, reported}
\end{center}

\begin{abstract}
PoLar (ICML 2026) showed that a frozen LLM admits many valid layer-execution
programs per input, and that a lightweight predictor can select one program that
is often shorter and more accurate than the standard forward pass. We ask a
different question: if several programs are valid, can we \emph{fuse} their outputs
rather than pick one? We introduce CROSS (Cross Router Of Selector System): K
candidate programs execute independently on the frozen backbone and a token-level
cross router fuses their hidden states with top-k weights before the LM head.
On a real 1B model, an in-sample fusion result (4.755) initially appeared to beat
the best single program (4.840), but 492K trainable parameters were fit on only
32--64 supervised tokens and both numbers used labels from the same windows; we do
\textbf{not} treat it as evidence of complementarity. On disjoint held-out windows
(\S\ref{sec:heldout}), the router-based reader reaches CE 3.061 (probability-level)
and 3.428 (hidden-level), below the identity program (3.847) and nearly matching
the labelled per-window oracle (3.056). However, the held-out fusion weights
collapse to a single program ([0,1,0,0] or [0,0,0,1]), and a larger held-out test
with paired bootstrap (14 windows, 210 tokens) gives fusion$-$P0 $= -0.075$ with
95\% CI $[-0.156, +0.012]$ --- crossing zero. We therefore withdraw both the
complementarity claim and the "programs beat identity" claim: the 3-window gain
was sampling noise. The one surviving positive result is a controlled one: with
\emph{fixed} mixing weights, probability-level fusion is below the Jensen bound
and below the best single program (\S\ref{sec:heldout}). We then
test the user-motivated idea of an execution-time selector that chooses which
programs to run \emph{before} execution, cutting compute. A distilled selector
perfectly mimics the full router (agree 1.00) and halves layer-executions
(93 $\to$ 46, $-$50.5\%), but CE rises by $+0.93$ nats: the six-step teacher router
under-valued the identity program, so the selector discarded the best candidate.
We report both results, the boundary they imply, and the full positioning relative
to PoLar, MoE layer-sharing, and our earlier negative result (BRICK). We claim
no novelty for programs, routers, MoE, or fusion individually; the contribution
is their intersection and its empirical capability boundary.
\end{abstract}

\section{Introduction}

This study continues our line of honest empirical boundary studies on small
language models. Amoeba established that verifier-based speculative decoding gives
a clean 2.015$\times$ on a full-attention 2B model but is architecture- and
harness-dependent. BRICK asked whether a 2D parallel-branch topology of a frozen
model enables dynamic activation; the answer was negative because a skipped
branch's information is genuinely absent and cannot be reconstructed by fusion
($+0.40$ nats irreducible loss).

CROSS asks the complementary question: what if instead of \emph{skipping}
information, we run several different computations and \emph{fuse} them? The
motivation comes from PoLar (Li et al., ICML 2026): pretrained layers can be
treated as a library of atomic functions, and a lightweight predictor generates
an input-specific execution program $\pi=(i_1,\ldots,i_K)$ that skips or repeats
layers. PoLar reports pass@k --- one of k programs is correct --- but always
returns a \emph{single} program's output. MoE work (MoRE, MoUE, UniPool, 2025--2026)
shares experts across layers \emph{within} the MoE layer. Nobody, to our knowledge,
fuses the outputs of multiple full program executions with a token-level router.
That intersection is CROSS. We stress that converting dense models into
MoE-like architectures is itself an active area (CMoE, ToMoE, ExpertWeaver);
our claim is limited to the program-as-expert variant, not to MoEfication in
general.

\section{Background}

\subsection{PoLar: program-of-layers}
A pretrained model of depth $D$ defines functions $f_0,\ldots,f_{D-1}$. A program
$\pi$ is a finite sequence of layer indices; its execution composes
$F_\pi=f_{i_K}\circ\cdots\circ f_{i_1}$. PoLar segments layers into contiguous
modules and assigns each \{skip, keep, repeat\}, predicting the segmentation and
operations with a 2.1M-parameter network. Findings: most valid programs are shorter
than the standard pass; recurrence helps more than skipping; combining both is
best; gains persist out of distribution. Its 24-layer Qwen1.5-MoE experiment reports
0.83$\times$--0.95$\times$ latency with $+5.7$--$9.4$ accuracy points.

\subsection{Why CROSS is not BRICK again}
In BRICK, skipping a branch at readout removed information: fusion could recover
the \emph{merging} loss but not the \emph{information} loss ($+0.40$ nats). CROSS
executes every candidate program fully, so no information is missing; fusion only
combines complementary complete representations. The cost is higher (K program
executions) --- CROSS is a capability study, not an accelerator.

\section{Method}

\subsection{Program pool}
We use K=4 fixed diverse programs on MiniCPM5-1B (24 layers): P0 identity (all
layers); P1 skip a middle segment; P2 repeat a middle segment; P3 skip a late
segment and repeat an early one. Each program keeps the final layer (readout
interface, from Amoeba/BRICK). This is a simplified PoLar-style pool; a learned
selector is introduced below.

\subsection{Cross router (fusion)}
Each program $j$ produces a final hidden state $h_j$. The cross router computes
token-level weights and fuses:
\[
  w(t)=\mathrm{softmax}\big(\mathrm{MLP}([h_1(t);\ldots;h_K(t)])\big),\quad
  w=\mathrm{topk\text{-}sparsify}(w,k{=}2),\quad
  h_{\mathrm{fused}}=\sum_j w_j(t)\,s_j\,h_j(t),
\]
where $s_j$ are learned scalar scales. Only next-token CE supervises the router.
Gumbel-Softmax and top-k sparsification prevent collapse to one program (BRICK's
lesson). Trainable parameters: 492K; the 2GB backbone is frozen.

\subsection{Execution-time selector (compute saving)}
The router's top-k only sparsifies fusion weights; all K programs still execute.
To actually save compute, a `ProgramSelector` MLP predicts program logits from a
cheap prelude hidden state (first 4 layers) \emph{before} any program runs, and only
the top-2 programs are executed. Hard top-k is non-differentiable, so the selector
is trained by distillation: its target is the normalized mean fusion weight of the
trained full router (teacher signal, not self-confidence --- the Amoeba v2 lesson).

\section{Experimental Setup}
Real MiniCPM5-1B (24 layers, hidden 1536), CPU-only (i7-1260P, 15.6GB, PyTorch
2.15 CPU, transformers 5.24.1), bf16 on a 2.36GB resident model. Small settings
(2--4 windows, seq 16, 6--12 steps) measure trends and boundaries, not final
accuracy. Two dtype bugs that crash on real bf16 models but never appear in fp32
tiny tests were found and fixed (same class as BRICK).

\section{Results}

\subsection{H1: fusion beats the best single program}
\begin{table}[h]
\centering
\begin{tabular}{lc}
\toprule
Metric & Value \\ \midrule
single-program CEs & 5.26 / 5.61 / 5.41 / 5.49 \\
best single (per-window oracle) & 4.840 \\
untrained fusion & 5.451 \\
\textbf{trained fusion (12 steps)} & \textbf{4.755} \\
\textbf{Delta vs best single} & $\mathbf{-0.085}$ \\
fusion weights & [0.17, 0.05, 0.25, 0.53] (no collapse) \\
trainable parameters & 492K \\
\bottomrule
\end{tabular}
\caption{CROSS H1 on real 1B: fusion of K programs vs single programs.}
\end{table}
On the training windows, the trained fusion beats every single program and the
per-window oracle, while untrained fusion (5.451) is worse than all singles. As
\S\ref{sec:heldout} shows, this in-sample result is fitting rather than
complementarity and the claim is withdrawn. We keep the table only as the
motivation for the held-out test.

\subsection{Held-out gatekeeper}
\label{sec:heldout}
Training used windows 0--2; evaluation used disjoint windows 3--5 (never seen in
training). Two fusion levels were tested: hidden ($h=\sum_j w_j h_j$ then LM head)
and probability ($p=\sum_j w_j\,\mathrm{softmax}(\mathrm{lm\_head}(\mathrm{norm}(h_j)))$,
which by Jensen's inequality has CE no worse than the weighted member average).

\begin{table}[h]
\centering
\begin{tabular}{lcccc}
\toprule
Split & P0 (identity) & oracle-min (labelled) & hidden fusion & prob fusion \\
\midrule
train (0--2) & 4.249 & 3.969 & 4.382 & 3.856 \\
\textbf{held-out (3--5)} & \textbf{3.847} & 3.056 & \textbf{3.428} & \textbf{3.061} \\
\bottomrule
\end{tabular}
\caption{Held-out evaluation: both fusion levels beat identity; probability fusion
almost matches the labelled oracle.}
\end{table}

Held-out fusion weights collapse to one program: hidden $[0,0,0,1]$ (P3),
probability $[0,1,0,0]$ (P1). The gain in that experiment therefore comes from
\emph{selecting a program better than identity}, not from fusing. On training
windows P0 was best (4.249); on held-out windows P1 (3.056) and P3 (3.428) beat
P0 (3.847) --- a ranking flip across splits. We initially attributed the
hidden-vs-probability gap (3.428 vs 3.061) to Jensen's inequality, but this was
\emph{confounded}: the two runs selected different programs, so both had one-hot
weights and the difference was 100\% selection, 0\% fusion level.

\paragraph{Larger held-out with paired bootstrap.} The 3-window held-out result is
underpowered (48 tokens). Re-running on 14 disjoint windows (210 tokens) with a
paired per-token bootstrap gives fusion 3.775, P0 3.850, oracle 3.240, and a
paired difference $-0.075$ with 95\% CI $[-0.156,+0.012]$. The CI crosses zero:
the apparent gain of selection over identity is \emph{not} statistically
significant and is withdrawn.

\paragraph{Clean fusion-level comparison.} To remove the confound we fix the
\emph{same} mixing weights for both fusion levels (uniform $[0.25]^4$ and the
in-sample weights), so the router is not selecting. Results on all 20 windows:

\begin{table}[h]
\centering
\begin{tabular}{lccc}
\toprule
fixed weights & hidden CE & probability CE & Jensen bound \\
\midrule
uniform $[0.25]^4$ & 3.827 & \textbf{3.719} & 3.966 \\
in-sample $[.17,.05,.25,.53]$ & 3.879 & \textbf{3.774} & 3.956 \\
\bottomrule
\end{tabular}
\caption{Fixed-weight fusion: probability fusion is strictly below the Jensen
weighted-member bound (3.966/3.956) and below the best single program (3.880).
Hidden fusion has no such guarantee.}
\end{table}

With identical weights, probability fusion (3.719) is below both the Jensen
bound (3.966) \emph{and} the best single program (3.880): a real fusion gain
with a theoretical guarantee. Hidden fusion (3.827) also beats the best single
but has no guarantee and is 0.11 worse than probability fusion. The supported
statement is twofold: (i) a frozen model admits programs better than identity on
unseen inputs and a router selects them; (ii) with fixed weights, probability
fusion reliably improves over both the member average (Jensen) and the best
individual program. The earlier ``complementarity via trained router'' claim is
withdrawn; this theory-backed fixed-weight result replaces it.

\subsection{Numeric discrepancies across experiments}
\begin{enumerate}[leftmargin=*]
  \item Trained fusion 4.755 (H1) vs 4.593 (selector Stage 1): the first used 12
        steps on 2 windows, the second 6 steps on 3 windows; the 0.16 gap is a
        training-budget/window-set difference, not a discrepancy. Both are
        superseded by the held-out numbers in \S\ref{sec:heldout}.
  \item P0 (identity) 5.26 in CROSS vs 3.79--4.04 in BRICK: different window sets,
        and BRICK's reference is a \emph{fused} both-branch path while CROSS's P0 is
        an unfused standard forward pass. Not comparable; we use CROSS's own P0
        throughout.
  \item BRICK's layers 5--21 $\approx 0$ ablation bypasses a \emph{single} layer with
        later layers compensating; CROSS's P1 skips a contiguous \emph{six-layer}
        segment (6--11), where the residual bypass compounds. Different
        interventions; the segment-level cost (+0.35 on train windows) is the
        relevant one for the program pool.
\end{enumerate}

\subsection{Selector: compute saving at a CE cost}
\begin{table}[h]
\centering
\begin{tabular}{lc}
\toprule
Metric & Value \\ \midrule
distill agreement (selector vs router) & 1.00 \\
selected programs & [1, 2] \\
full CROSS CE (93 layer-executions) & 4.593 \\
select-before-execute CE (46 layer-executions) & 5.528 \\
CE delta & $+0.934$ \\
compute saving & 50.5\% \\
\bottomrule
\end{tabular}
\caption{Execution top-k: half the compute, higher CE (negative result).}
\end{table}
The selector perfectly learns the teacher router's choice and halves compute, but
CE rises by $+0.28$--$0.93$ across two runs with 6- vs 10-step teachers: the
under-trained teacher's program ranking is unstable (one run selected [1,2] and
dropped identity; another selected [0,1]), and forcing P0 in improves CE by 0.09
but does not remove the gap. Execution top-k is thus highly sensitive to teacher
training --- a boundary result, not a final verdict. We record it honestly.

\section{Negative Results and Boundaries}
\begin{enumerate}[leftmargin=*]
  \item \textbf{Untrained fusion hurts} (5.451 vs best single 4.840): naive averaging
        of program outputs is worse than every program.
  \item \textbf{Execution top-k loses capability} ($+0.93$) when the teacher router
        is under-trained. Compute saving is real (50.5\%), but program choice matters.
  \item \textbf{Tiny models hide dtype and routing bugs}: fp32 toy tests passed while
        bf16 real-model runs crashed twice; also the router collapsed on tiny but
        not the same way on real data.
  \item \textbf{Small sample}: 2--4 windows and 6--12 steps give trends; full
        convergence and 79-question evaluation require a GPU.
  \item \textbf{Cost}: K program executions make CROSS a capability study; it does
        not accelerate inference.
\end{enumerate}

\section{Related Work (positioning, not novelty claims)}
\begin{itemize}[leftmargin=*]
  \item \textbf{Dense-to-MoE conversion (MoEfication)}: CMoE (arXiv:2502.04416)
        converts a dense LLM to MoE without training; ToMoE (ICML 2026) uncovers
        experts by structural pruning; ExpertWeaver (arXiv:2602.15521) builds a
        training-free router from GLU gating centroids; DOT-MoE (2026) partitions
        FFN neurons into experts. \textbf{The dense-to-MoE idea itself is not novel
        and we do not claim it.} Their experts are \emph{intra-layer substructures}
        (FFN neurons, GLU groups, pruned blocks). CROSS is distinct in treating an
        entire \emph{layer-execution program} (skip/repeat layer segments) as the
        expert and fusing full program outputs. We claim only this specific
        program-as-expert intersection, not MoEfication in general.
  \item \textbf{PoLar} (ICML 2026): single-program skip/repeat selection on frozen
        layers. CROSS fuses multiple program outputs; we do not claim the program
        representation is novel.
  \item \textbf{MoE cross-layer sharing}: MoRE, ReXMoE, MoUE, UniPool (2025--2026)
        share experts within layers; CROSS fuses program-level outputs.
  \item \textbf{State fusion}: Mixture of States (2025) mixes layer hidden states
        for multimodal generation; Attentive Multi-Layer Fusion (ICML 2026) fuses
        vision-transformer layers. Neither is LLM program fusion.
  \item \textbf{Earlier work}: LayerSkip, BitSkip (early exit); Amoeba (capability
        criterion, speculative decoding); BRICK (2D topology negative result).
\end{itemize}

\section{Conclusion}
CROSS asked whether MoE-style fusion over multiple program outputs can improve a
frozen small model beyond identity. The in-sample result (4.755) was rejected as
fitting. On held-out windows, a trained router-based reader reaches CE 3.061
(probability-level), below identity (3.847) and near the labelled oracle (3.056),
but the weights collapse to one program. A larger held-out test with paired
bootstrap reverses the interpretation: fusion$-$P0 $= -0.075$, 95\% CI
$[-0.156,+0.012]$, crossing zero. We therefore withdraw both complementarity and
the "programs beat identity" claim --- the original gain was sampling noise. An
execution-time selector saves about 50\% of compute but raises CE by $+0.28$--$0.93$
and is unstable across teacher budgets. The single supported conclusion is a
controlled one: with fixed mixing weights, probability-level fusion is below the
Jensen weighted-member bound (3.966) and below the best single program (3.880),
while hidden-level fusion has no such guarantee. We release all code and logs.
The contribution is an empirical boundary map, not a novelty claim.

\end{document}
